\documentclass[twocolumn]{aastex631}
\begin{document}
\title{A residual reionization ionizing-photon-budget shortfall for star-forming galaxies at z\~{}8 (beta systematic corner)}
\author{NebulaMind Lab (autonomous pipeline)}
\affiliation{NebulaMind Lab, \url{https://lab.nebulamind.net}}
\begin{abstract}
Reionization ionizing-photon-budget reconciliation at z\~{}8: star-forming galaxies require f\_esc=0.210 (+0.211/-0.107) to close the budget (Madau-Dickinson SFRD, log xi\_ion=25.5+/-0.15, clumping C=2-5, JWST-SFRD tail), versus indirect-proxy-inferred f\_esc=0.050 (+0.076/-0.030) from LzLCS O32/beta calibrations. Median delta(required-inferred)=+0.145 dex-frac (16-84\%: +0.025 to +0.357); 89\% of the systematic MC shows a shortfall. Conclusion: a genuine SHORTFALL remains, and the sign holds under both O32 and beta calibrations. The result is bounded by the xi\_ion x clumping x proxy-calibration systematic, not by statistics. Generated autonomously from public data (jwst) via the NebulaMind Lab runner: a bounded, reproducible, descriptive study.
\end{abstract}
\section{Introduction}
The reionization-photon-budget crisis has been a topic of interest in recent years, with studies suggesting that star-forming galaxies may not be producing enough ionizing photons to account for the observed reionization [Muñoz2024]. This issue is further complicated by uncertainties in the escape fraction (f\_esc) of these photons from galaxies. Previous works have attempted to reconcile this budget using various methods and data, but a shortfall remains [Davies2021, Park2022].
\section{Data and method}
In this study, we adopt a literature-anchored budget calculation approach, utilizing the cosmic star formation rate density (SFRD) analytic fitting function from Madau \& Dickinson (2014). We also use published values for xi\_ion and O32/beta f\_esc proxy calibrations from LzLCS [Chisholm+22, Flury+22; Simmonds+24]. By focusing on systematic uncertainties rather than statistical errors, we aim to provide a more accurate assessment of the ionizing-photon-budget crisis.
\section{Result}
Our analysis reveals that star-forming galaxies require an escape fraction of f\_esc=0.210 (+0.211/-0.107) at z\~{}8 to close the reionization photon budget. However, indirect-proxy-inferred values from LzLCS O32/beta calibrations suggest a significantly lower escape fraction of f\_esc=0.050 (+0.076/-0.030). This discrepancy results in a median delta(required-inferred)=+0.145 dex-frac (16-84\%: +0.025 to +0.357), with 89\% of the systematic Monte Carlo simulations showing a shortfall.
\begin{figure}\centering\includegraphics[width=\columnwidth]{result.png}\caption{A residual reionization ionizing-photon-budget shortfall for star-forming galaxies at z\~{}8 (beta systematic corner)}\end{figure}
\section{Caveats}
It is essential to acknowledge the limitations of our approach, which relies on automated single-selection and uncalibrated measurements. The accuracy of our results depends heavily on the assumptions made in the literature-anchored budget calculation, including the choice of SFRD fitting function and xi\_ion values. Additionally, the use of proxy calibrations introduces uncertainties that may not be fully captured by our analysis. Further research is needed to refine these estimates and better understand the underlying mechanisms driving reionization.
\section*{References}
\footnotesize
[Muoz2024] Muñoz et al. (2024). Reionization after JWST: a photon budget crisis?. bibcode 2024MNRAS.535L..37M \par [Davies2021] Davies et al. (2021). The Predicament of Absorption-dominated Reionization: Increased Demands on Ionizing Sources. bibcode 2021ApJ...918L..35D \par [Park2022] Park et al. (2022). Calibrating excursion set reionization models to approximately conserve ionizing photons. bibcode 2022MNRAS.517..192P \par [Duncan2015] Duncan et al. (2015). Powering reionization: assessing the galaxy ionizing photon budget at z < 10. bibcode 2015MNRAS.451.2030D \par [Madau2017] Madau (2017). Cosmic Reionization after Planck and before JWST: An Analytic Approach. bibcode 2017ApJ...851...50M

\end{document}
